\section{Adjoint positivity of an extreme Hodge line}
\label{period:section}

The relative canonical bundle calculation gives a line in the highest
nonzero Hodge filtration step of one complex summand.  Its curvature
is semipositive, but this alone does not compare it with a canonical
bundle.  We construct a projective quotient on which the line descends
and prove that adding a positive multiple of it makes the canonical
class big.

Here an \emph{extreme Hodge line} means the rank-one highest nonzero
filtration step of a complex Hodge direct summand.  An integral
structure is a lattice in the ambient real local system.  The chosen
summand need not be rational, and the ambient real polarization need
not be rational either.

The \emph{parabolic extension} is obtained by local power substitutions
that make boundary monodromy unipotent, followed by Schmid extension
and descent with rational weights.  The monodromy theorem gives
rational weights in the integral case.  When the local monodromy has
a finite part, these weights are essential.  We first resolve any
compactification whose boundary is not simple normal crossing.
Unipotent extension commutes with pullback, so the same is true of
the parabolic rational line after the power substitutions; see
\cite{Schmid,CKS} and \cite[Lemma~5.6]{BFMT}.  We use version~2
of \cite{BFMT} throughout.


We prove the adjoint-positivity theorem stated as
\Cref{period:adjoint}. The line in that statement is allowed to be a
complex character line inside an integral ambient variation.

There are two maps in the proof.  In a flat trivialization the
\emph{line map} records only the chosen highest line; the \emph{period
map} records all Hodge filtration steps.  Their ranks can differ.
Curvature identifies the numerical dimension of the line with the
rank of the first map.  We use the full periods of suitable rational
adjoint factors to construct the projective quotient.  On that
quotient, logarithmic general type gives a big canonical class with
boundary.  The remaining argument proves that the line loses numerical
rank on each marked boundary component, so a section count removes
the boundary after adding a multiple of the line.

\subsection{Curvature, numerical dimension, and the boundary}

We first relate the curvature on the open set to intersection numbers
on the compactification.  For a nef rational line bundle $A$ on a compact
Kähler $n$-fold, we use
\[
 \nu(A)=\max\{k:c_1(A)^k\cdot[\omega]^{n-k}>0\},
\]
where $\omega$ is any Kähler form.  On a projective manifold this is
the usual numerical dimension.

\begin{lemma}[Curvature and numerical dimension]
\label{period:numerics}
Let $T$ be a smooth compact Kähler manifold, let $E\subset T$ be a
simple normal crossing divisor, and let $A$ be the parabolic extension
of the rank-one highest step of a polarizable pure complex variation on
$T\setminus E$ with quasi-unipotent local monodromy.  Then $A$ is nef.
Its numerical dimension equals the generic rank of the map that
remembers this line in a flat trivialization.

The corresponding assertion for the Griffiths line, the tensor product
of the determinants of the Hodge filtration steps, uses the full period
map.  In particular, its rational extension is big if the full period
map is generically immersive.
\end{lemma}

The curvature-and-volume argument is closely related to the proofs of
\cite[Lemmas~6.15 and~6.17]{BakkerBrunebarbeTsimerman}; here we need
all mixed intersection numbers, since the chosen line need not be big.

\begin{proof}
We may check the assertions after a finite cover and resolution on
which local monodromy is unipotent.  This does not require an integral
lattice.  Indeed, $\pi_1(T\setminus E)$ is finitely generated.  Put
the entries of finitely many monodromy generators and their inverses
in a finitely generated $\Z$-algebra $R\subset\C$.  Adjoin the
finite group of roots of unity generated by the eigenvalues of the
finitely many boundary monodromies, and invert every difference
$\zeta-1$ for a nonidentity element of this group.  Reduction at a
maximal ideal with finite residue field gives a finite-index normal
congruence subgroup.  Every element of a local boundary monodromy
group that lies in this subgroup is unipotent: its eigenvalues
belong to that finite group and reduce to one, so the inverted
differences force them to equal one.  This also applies to products
of commuting boundary monodromies at crossings.  The corresponding finite cover extends
over the compactification, and a Kähler resolution gives the desired
unipotent model; see \cite[Lemma~2.21]{VilladsenHodgeKahler}.
These covers preserve nefness, numerical dimension, and bigness of
the rational line.  Let $h$ be its Hodge metric and $\Theta$ its curvature
on $T\setminus E$, normalized to represent $c_1(A)$.  Griffiths'
curvature formula gives $\Theta\geq0$; its rank is the rank of the
line map, since its value on a tangent vector is the squared norm of
the second fundamental form of the highest step.  The same formula
for the Griffiths line detects the differential of the full period
map.

The Hodge norm estimates in extending frames give, in boundary
coordinates $t_1,\ldots,t_k$, weights whose absolute values are bounded
by
\begin{equation}
\label{period:loglog}
              C\left(1+\sum_{j=1}^k
                   \log(-\log|t_j|)\right).
\end{equation}
The semipositive metric therefore extends as a positive current in
$c_1(A)$ with zero Lelong numbers.  Regularization of positive
currents then gives nefness; see
\cite[Corollary~6.4]{DemaillyRegularization}.
This extends the nefness argument of \cite[Lemma~2.16(1)]{BFMT}
from a proper log smooth algebraic space to the present compact
Kähler setting.

To identify numerical dimension, we must also show that no
intersection mass is lost at $E$.  Write $\alpha=c_1(A)$ and fix
$\varepsilon>0$.  Since $\alpha$ is nef, $\alpha+\varepsilon[\omega]$
has a smooth Kähler representative
$\theta_\varepsilon$.  Write
\[
 \Theta+\varepsilon\omega=\theta_\varepsilon+\ddc u,
 \qquad \ddc=\frac{\sqrt{-1}}{2\pi}\partial\bar\partial.
\]
The global potential $u$ satisfies the local bounds
\eqref{period:loglog}. Put $n=\dim T$. If $E$ is empty, the potential
is smooth and there is no lost mass. Otherwise write
$E=E_1+\cdots+E_N$ and choose smooth divisor metrics, with defining
sections $s_j$, such that
\[
 u_j=-\log|s_j|^2,\qquad
 \ddc(-u_j)=[E_j]-\eta_j,
\]
where the $\eta_j$ are smooth curvature forms. Choose $c>0$ with
$c\eta_j\leq\theta_\varepsilon$ for all $j$, and rescale the metrics
so that $cu_j\geq1$. Then $\phi_j=-cu_j$ is
$\theta_\varepsilon$-plurisubharmonic and $\phi_j\leq-1$.
For any fixed $0<b<1$, the attenuation theorem
\cite[Example~2.14]{GuedjZeriahiEnergy} gives
\[
 \psi_j=-(-\phi_j)^b=-c^bu_j^b
\]
with full Monge--Amp\`ere mass in the K\"ahler class
$[\theta_\varepsilon]$. The full-mass class is convex and contains
the zero potential \cite[Proposition~1.6]{GuedjZeriahiEnergy}.
Thus, for $0<\delta\leq c^b/N$, the potential
\begin{equation}\label{period:barrier}
 v=-\delta\sum_{j=1}^N u_j^b
   =\sum_{j=1}^N\frac{\delta}{c^b}\psi_j
      +\left(1-\frac{N\delta}{c^b}\right)0
\end{equation}
is $\theta_\varepsilon$-plurisubharmonic and has full mass. A positive power grows
faster than a logarithm, so \eqref{period:loglog} implies
$v-C\leq u$ globally for some constant $C$.
Full mass is preserved on passing to less singular potentials
\cite[Propositions~2.11(i) and~2.14(i)]{BEGZ}. Hence $u$ has full mass.
The non-pluripolar measure gives no mass to the analytic set $E$ and
agrees with the smooth wedge product on its complement. Consequently
\[
 \int_{T\setminus E}(\Theta+\varepsilon\omega)^n
                  =(\alpha+\varepsilon[\omega])^n.
\]
Expand both sides as polynomials in $\varepsilon$.  Their coefficients
are finite, since all summands on the left are nonnegative.  Equality
for all $\varepsilon>0$ gives
\begin{equation}
\label{period:masses}
 \int_{T\setminus E}\Theta^k\wedge\omega^{n-k}
                         =\alpha^k\cdot[\omega]^{n-k}
                         \qquad(0\leq k\leq n).
\end{equation}
If the generic curvature rank is $r$, the left side is positive for
$k\leq r$ and zero for $k>r$.  This proves the numerical assertion.
For an immersive Griffiths line the top integral is positive, and
the volume criterion for a positive current on a compact Kähler
manifold gives bigness; see \cite{BoucksomVolume}.
\end{proof}

\begin{lemma}[Restriction to a boundary component]
\label{period:boundary}
In the unipotent case of \Cref{period:numerics}, let $E_i$ be a smooth
boundary component and $N_i$ its monodromy logarithm.  On the open
stratum of $E_i$, the limiting highest line occurs in a unique weight
grade of the limiting mixed variation.  The restriction $A|_{E_i}$
is the Schmid extension, across the remaining boundary of $E_i$, of
that graded Hodge line.  Its numerical dimension is therefore the
generic rank of the graded line map.
\end{lemma}

\begin{proof}
The extended filtration and $W(N_i)$ give the limiting mixed
variation along the open stratum.  Since the highest step has rank
one, exactly one weight grade contains it.  Changes of the normal
parameter act by $\exp(cN_i)$ and act trivially on the associated
graded.  At a further crossing the relative monodromy filtration
identifies the extension of this graded variation with the graded
of the original extending filtration.  Thus the identification holds
on the entire component.
For log smooth algebraic spaces it is stated in
\cite[\S2.5.3 and Lemma~2.16(3)]{BFMT}; its local proof uses the
multivariable nilpotent orbit theorem \cite{CKS} and applies here
as well.  The filtrations and strictness also restrict to a complex
summand.  Apply \Cref{period:numerics} on $E_i$.
\end{proof}

\subsection{Rational adjoint periods and descent of the line}

The period-quotient construction requires discrete monodromy.  A
projection to one complex factor need not retain it.  We therefore
construct a rationally polarized adjoint variation from the ambient
integral local system and keep whole rational factors.  A tensor
construction then recovers a positive power of the original line,
including its scalar monodromy.

\begin{lemma}[Rational adjoint reduction]
\label{period:adjoint-reduction}
In the setting of \Cref{period:adjoint}, let $G$ be the identity
component of the rational Zariski closure of monodromy.  Then $G$ is
semisimple.  Its adjoint local system is a rationally polarized
integral variation of weight zero.

There is a monodromy-invariant collection of rational simple adjoint
factors, each containing a complex factor acting nontrivially on the
highest-line orbit, with the following properties.  The resulting
adjoint variation $\mathbb A$ has discrete integral monodromy.  After a
positive power and up to a constant one-dimensional complex Hodge
shift, the original line is the highest line of a complex
subvariation of a tensor construction on $\mathbb A$.  On a finite
cover preserving the factors, that line splits into the tensor
product of the highest lines belonging to the individual complex
factors.  All these highest steps have rank one.
\end{lemma}

\begin{proof}
Semisimplicity and the theorem of the fixed part hold for polarizable
real and complex variations on a Zariski open subset of a compact
Kähler manifold.  We may use a Kähler modification of $B$ to apply
them; see \cite[Lemma~2.3 and its proof]{PopaSchnellKahler} and
\cite[Theorem~3 and Remark~4]{BakkerMatsushita}.  Thus the algebraic
monodromy group is reductive.  Pass temporarily to a finite cover
on which it is connected.  The irreducible complex constituents of
the local system carry polarizable complex variations.  They can be
defined over a number field: the representation is rational and its
algebraic monodromy is reductive.  The determinant of each constituent
is unitary, as are all its algebraic conjugates, which are again
constituents of the rational representation.  Its monodromy values
are algebraic units by the integral lattice.  Kronecker's theorem
makes these values roots of unity in a fixed number field, hence
makes the determinant character finite.  It is consequently trivial
on $G$.  This is Deligne's finite-determinant argument; its classical
algebraic formulation is \cite[Corollary~4.2.8(iii)(b)]{DeligneHodgeII}.
The connected center of $G$ acts by scalar characters on
the irreducible constituents; their determinants detect these
characters.  Since the original representation is faithful, this
connected center is trivial.  Therefore $G$ is semisimple.

We next equip the monodromy Lie algebra with a rational polarization.
In every tensor construction,
the subspace of $G$-invariant tensors is the fixed part and has a
constant Hodge structure.  In a flat trivialization the Hodge circles
$h_x$ and $h_{x_0}$ therefore induce exactly the same linear operator
on each such subspace at each parameter $z\in S^1$.  Reductivity
identifies $G$ with the pointwise stabilizer of all its invariant
tensors.  The circles preserve these subspaces, so normalize $G$,
and the equality of their actions gives
\[
                    h_x(z)h_{x_0}(z)^{-1}\in G(\R).
\]
Here is a direct local conjugacy argument.  Along a smooth path in
the base, write $a_t=\operatorname{Ad}\circ h_t:S^1\to G^{\rm ad}(\R)$
and $u_t(z)=(\partial_t a_t(z))a_t(z)^{-1}$.  Differentiating the
homomorphism identity gives
\[
 u_t(zz')=u_t(z)+\operatorname{Ad}(a_t(z))u_t(z').
\]
Normalized Haar averaging, with $X_t=\int_{S^1}u_t(z')\,dz'$, yields
$u_t(z)=X_t-\operatorname{Ad}(a_t(z))X_t$.  The solution of
$g_t'g_t^{-1}=X_t$, $g_0=1$, therefore gives
$a_t=g_ta_0g_t^{-1}$.  Lift this path through the central isogeny
$G(\R)^+\to G^{\rm ad}(\R)^+$.  The ratio of $h_t$ and
$g_th_0g_t^{-1}$ lies in $G(\R)$ by the invariant-tensor equality,
and centralizes $G$ by the equality of their adjoint actions.
It lies in the finite center of $G$; continuity in the circle
parameter makes it the identity.  Thus $h_t=g_th_0g_t^{-1}$,
proving that the lifted period map lies in a single
$G(\R)^+$-orbit.

It follows that $\mathfrak g_{\Q}=\operatorname{Lie}(G)$ is a Hodge
sub-local system of $\End(\mathbb V_{\Q})$, with its weight-zero
grading.  The Weil involution $\theta$ preserves $\mathfrak g_{\R}$
and is the restriction of the Cartan involution of the real
polarization isometry group.  To check the restriction, use the
positive Hodge metric: on its Lie algebra $\theta(X)=-X^*$.
If $\mathfrak g_{\R}=\mathfrak k\oplus\mathfrak p$ is its
eigenspace decomposition, then $\mathfrak k\oplus i\mathfrak p$
is a compact real form of $\mathfrak g_{\C}$, represented by
skew-Hermitian endomorphisms.  Thus
$-B_{\mathfrak g}(X,\theta X)$ is positive definite, proving that
$\theta$ is Cartan on $\mathfrak g_{\R}$.  The appropriately
signed Killing form therefore polarizes this weight-zero variation.
The form is rational.  Moreover
\[
          \mathfrak g_{\Q}\cap\End(\mathbb V_{\Z})
\]
is a monodromy-invariant lattice.  This gives the rationally polarized
integral adjoint variation.

Let $\ell$ be the highest line at a reference point of the chosen
complex summand.  The non-lowering parabolic of the Hodge grading
preserves $\ell$, so a Borel subgroup preserves it.  Its closed orbit
\[
                       J=G_{\C}\ell
\]
is a flag variety, a product of flag varieties for the complex simple
factors.  The span $W$ of this orbit is irreducible.  Indeed, in a
semisimple decomposition of the representation, the nonzero
components of a Borel-eigenline all have the same highest weight;
their span is a single diagonal copy of that irreducible
representation.  The line is its unique highest Hodge step.  Indeed, at every point
the Hodge circle normalizes $G$ and preserves the highest line there;
it consequently preserves the span of its $G_{\C}$-orbit.  Thus $W$
is a Hodge subvariation.  The lifted line map takes values in $J$,
and the original line is the pullback of the tautological line there,
equivariantly for monodromy.

Call a complex simple factor \emph{visible} if it acts nontrivially
on $J$.  Retain every rational simple adjoint factor containing a
visible factor.
This collection is invariant under the possibly disconnected
monodromy group.  On its Lie algebra we obtain the variation
$\mathbb A$; its monodromy is contained in the automorphisms of a
lattice and is discrete.  Notice that an entire rational factor is
kept.  Projection to only one real factor would not ensure
discreteness.

We now recover the line from this adjoint system.  Choose a positive
integer $N$ for which the highest weight $N\lambda$ of $W$ belongs
to the root lattice.  The Cartan component
$W_N\subset W^{\otimes N}$ is generated by the orbit of
$\ell^{\otimes N}$.  It factors through the connected adjoint group,
and its highest Hodge line is $\ell^{\otimes N}$.  The same
orbit-span argument just used shows that $W_N$ is a Hodge
subvariation.

We must also account for disconnected monodromy before applying
tensor generation.  Deck transformations take the lifted line to
another point of the same connected orbit, so they preserve $W$ and
its Cartan components.  The kernel of the action on the retained
adjoint factors centralizes the $G$-action on $W$, and hence acts
scalarly by irreducibility.  Its connected part is semisimple and
has no nontrivial scalar character.  Its component group is finite.
Increasing $N$ therefore kills the entire kernel action on $W_N$.
Consequently $W_N$ is a representation of the full projected
algebraic monodromy group, not just of its identity component.
This step rules out a discarded nontorsion scalar local system.

The adjoint representation of that projected group is faithful and
self-dual.  By tensor generation and complete reducibility,
$W_N$ embeds as a flat direct summand in a tensor construction
$\mathbb T(\mathbb A)_{\C}$; see
\cite[Theorem~4.14]{MilneAlgebraicGroups}.  To make the embedding
compatible with Hodge structures, apply the theorem of the fixed part
to
\[
 H^0\bigl(B^\circ,
       \mathcal Hom(W_N,\mathbb T(\mathbb A)_{\C})\bigr).
\]
This nonzero space of flat maps has a constant complex Hodge
structure.  A nonzero Hodge-homogeneous component of any nonzero
flat map is still flat.  Its kernel is $G$-invariant, so the map is
injective because $W_N$ is irreducible.  A constant one-dimensional
complex Hodge shift makes it a morphism of variations; that shift
changes neither the underlying holomorphic line nor its boundary
extension.  This is the constituent and fixed-part argument of
\cite[Theorem~3]{BakkerMatsushita}.

The complex summand carries a flat Hermitian polarization.  In a
real tensor variation of weight $w$ polarized by $Q$, it is induced
by $i^wQ_{\C}(v,\bar u)$, up to the conventional overall sign.
Its restriction to a Hodge subvariation is nondegenerate because
it is definite with the prescribed sign on each Hodge component.
There is no need for the real bilinear form itself to restrict
nondegenerately to the chosen complex summand: it may pair that
summand with its conjugate.

Finally, on a factor-preserving finite cover an irreducible
representation is a tensor product over the simple factors.  The
dimension of its highest step is the product of the corresponding
dimensions.  Since that dimension is one, each factor has a unique
rank-one highest step.  This gives the asserted factor lines.  Their
extensions are nef and functorial by
\Cref{period:numerics,period:boundary}.
\end{proof}

\begin{lemma}[Generic stabilizers of adjoint periods]
\label{period:stabilizer}
Let $G$ be a connected semisimple adjoint group over $\Q$, and let
$\mathbb A_{\Q}$ be a rationally polarized integral variation whose
fiber is $\mathfrak g_{\Q}=\operatorname{Lie}(G)$ and whose connected
algebraic monodromy is $G$, acting by its adjoint representation.
No nonidentity rational Lie algebra automorphism of
$\mathfrak g_{\Q}$ fixes all generic lifted periods of $\mathbb A$.
The automorphism need not belong to the monodromy group.
\end{lemma}

\begin{proof}
Suppose $\gamma$ fixes those periods.  By equivariance, every
monodromy conjugate of $\gamma$ is everywhere of Hodge type $(0,0)$
in the endomorphism variation.  Their rational span is therefore a
rational subvariation of type $(0,0)$.  Its induced polarization is
positive definite.  Intersection with the endomorphism lattice is
a full monodromy-invariant lattice: clear one denominator of
$\gamma$; integral monodromy conjugation preserves that denominator.
A positive definite integral monodromy group is finite.  Consequently
the connected group $G$ fixes $\gamma$, so $\gamma$ centralizes every
inner automorphism of $\mathfrak g$.  The centralizer of the inner
group in $\operatorname{Aut}(\mathfrak g)$ is trivial: on each simple
ideal commutation with all adjoint operators makes an endomorphism
scalar, and compatibility with the Lie bracket makes this scalar
one; permutations of distinct ideals do not commute with their
separate inner actions.  Thus $\gamma=1$.
\end{proof}

\subsection{A projective quotient for the retained periods}

We now quotient by connected fibers of the retained full period map.
The preceding reduction supplies both discrete integral monodromy
and a tensor description of a power of the line.  The first property
allows us to compactify the quotient; the second will descend the
line with its exact parabolic extension.  No bigness assertion about
$K_S+aP$ is used in this construction.

\begin{proposition}[The projective period quotient]
\label{period:quotient}
In \Cref{period:adjoint}, after modifications there is a map
$p:B'\to S$ with connected fibers to a smooth projective variety
and a nef rational line $P$ with $\tau^*L=p^*P$.  On a dense open
subset $S^\circ$, the retained adjoint variation of
\Cref{period:adjoint-reduction} descends, has the same connected
algebraic monodromy, and has a generically immersive period map.
The line $P$ is its parabolically extended highest-weight line,
up to the positive scaling already described.  If the retained
collection is empty, $L$ is rationally trivial and $S$ may be a
point.
\end{proposition}

The construction uses the compact-period-quotient approach of
\cite[Theorems~7.6--7.7]{BakkerBrunebarbeTsimerman}, whose
compactification input goes back to
\cite[Proposition~III and Remark~III-C]{SommesePeriodMapping}.
We use the precise class-$\mathcal C$ formulation of
\cite[Theorem~2.19]{VilladsenHodgeKahler} below.

\begin{proof}
Choose a compact K\"ahler modification of the class-$\mathcal C$
source. Then pass to a finite level at which the projected monodromy is
neat and preserves the factors.  Finite covers of the open locus
extend to finite ramified covers of compactifications and then to
Kähler resolutions; see \cite[Lemma~2.21]{VilladsenHodgeKahler}.
Resolve the boundary to simple normal crossings.  Extend the period
map over all finite-local-monodromy strata.  At neat level their
monodromy is trivial and the nilpotent orbit theorem gives this
extension.

The extended period map on this locus is proper; see
\cite[Proposition~9.11 and the proof of Theorem~9.6]{GriffithsPeriodsIII}.
Indeed, if a sequence tending to the boundary has period images in
a compact subset of the quotient, translate period lifts into a
compact set.
The punctured-disk Schwarz estimate makes the displacement of each
vanishing-coordinate monodromy tend to zero.  The discrete group
acts properly, since its isotropy in the real period domain is
compact.  The local monodromies must therefore be trivial at neat
level, so such a sequence cannot escape the extended locus.

Take the Stein factorization of this proper map.  Its first map has
compact connected fibers and its base is finite over the period
image.  It extends after modifications to a map of compact
class-$\mathcal C$ spaces with Kähler resolutions by
\cite[Theorem~2.19]{VilladsenHodgeKahler} (version~1).
That theorem applies to proper connected-fiber maps from Zariski
open subsets of compact class-$\mathcal C$ spaces.  In terms of
cycles, one first flattens the proper map, sends its irreducible general fibers
to a component of the Douady space of the compact source, and takes
the closure of that family.  The component is compact and in class
$\mathcal C$.  Near a general fiber the map to the cycle space is an
open embedding, since an irreducible compact cycle of the same
dimension in a sufficiently small neighborhood maps to a compact
analytic subset of a small Stein neighborhood in the base, hence
to a point, and must be the whole irreducible fiber.
The incidence map is consequently a modification of the source.
This produces the required compact quotient and also shows its
uniqueness as the quotient by the connected general period fibers.

On a dense open subset of this quotient the period and its integral
variation descend.  At neat level a constant period has trivial
stabilizer; local slices to the connected-fiber map give the descent
data.  After shrinking to a proper smooth fibration, the map on
fundamental groups is surjective, so the descended variation has
the same connected algebraic monodromy.  Its period differential is
generically injective by construction.  On a Kähler compactification
its Griffiths line is big by \Cref{period:numerics}; equivalently,
one can apply \cite[Corollary~1.3]{BrunebarbeCadorel}.  The quotient
therefore has a smooth projective model.  This is also the argument
in \cite[proof of Corollary~2.22]{VilladsenHodgeKahler}.

The finite deck group acts on the connected-period-fiber quotient.
Take its finite quotient and resolve.  This again has a smooth
projective model.  A general fiber before this finite quotient is
the union of connected level fibers indexed by one deck-group orbit.
The group acts transitively on these fibers, so the quotient is the
continuous image of any one of them and is connected.  Thus the
induced map from the original source has connected general fibers,
hence connected fibers after its Stein factorization.  The latter
does not change its function field.
All maps can be resolved using modifications of the original
source and of this projective model.

There is an actual adjoint variation on a dense Zariski open subset
of the quotient before changing level.  Write $\Gamma'$ for the
normal level subgroup of $\Gamma$.  The induced action of
$\Gamma/\Gamma'$ on the level quotient is faithful.  Indeed, if
$\delta\in\Gamma$ acts trivially there, discreteness shows locally
that its action on lifted periods agrees with a fixed
$\eta\in\Gamma'$.  Then $\eta^{-1}\delta$ fixes a period germ,
hence the whole generic lifted period image by analytic
continuation.  By \Cref{period:stabilizer} it is the identity,
so $\delta\in\Gamma'$.  Resolve the finite group action
equivariantly and remove its proper algebraic fixed loci and the
branch locus.  On the resulting dense Zariski open subset the level
quotient is free and unramified.  Its equivariant integral adjoint
local system and Hodge filtration descend there.  In particular,
this descent uses an open subset, not merely very general points.

The highest-weight construction in
\Cref{period:adjoint-reduction} gives a line on that open subset
whose pullback is a fixed positive power of $L$.  Take its
parabolic extension on a log smooth projective compactification
$S$ and divide by that power to obtain $P$.  Nefness follows from
\Cref{period:numerics}.  Pullback of the unipotent extensions,
followed by rational descent, gives
$\tau^*L=p^*P$ on the compact models, including the boundary.
Equivalently, this can be checked at level and then descended by
the norm: a line whose finite pullback is trivial is torsion in
the rational Picard group.  The identity therefore includes the
boundary and persists under further modifications.

If no factor is retained, the orbit is a point and the only possible
scalar monodromy is the finite scalar group treated in
\Cref{period:adjoint-reduction}.  Its parabolic line is rationally
trivial, as asserted.
\end{proof}

\subsection{Loss of numerical rank on the marked boundary}

Fix the projective quotient of \Cref{period:quotient} and a log smooth
compactification $(S,E)$ of its descended variation.  Let $D\leq E$
be the reduced sum of components with nonidentity local monodromy
on the retained adjoint system, including finite monodromy.
Logarithmic general type will give bigness of $K_S+D$.  The next
lemma is what permits us to remove $D$.

\begin{lemma}[Strict loss of rank at the marked boundary]
\label{period:rankloss}
For every component $D_i$ of $D$,
\begin{equation}
\label{period:strictdrop}
                       \nu(P|_{D_i})<\nu(P).
\end{equation}
\end{lemma}

\begin{proof}
Write $n=\dim S$ and $r=\nu(P)$, the generic rank of the line map
$j$ by \Cref{period:numerics}. The full retained period map has rank
$n$; the argument does not require $r=n$.
Work first on a finite Galois cover of the
open locus with connected algebraic monodromy and no factor
permutations.  Its connected adjoint group is denoted by $G$.

\smallskip\noindent
\emph{The Ax--Schanuel decomposition.}
Choose a Hodge-generic point, also general for the ranks of the
factor period maps and their joint differentials.  There are only
finitely many rational normal-factor decompositions of $G$, so
these requirements can be imposed simultaneously.  Let $T$ be a
local fiber germ of $j$ there, and $Z$ its irreducible Zariski
closure in the source of the variation.  Thus $\dim T=n-r$.
Move within the germ to a smooth point of $Z$ and use a smooth
open subset of a resolution when applying Hodge theory.

The restricted variation on $Z$ has the same generic Mumford--Tate
group as the ambient variation, since $Z$ contains the chosen
Hodge-generic point.  The monodromy normality theorem therefore
makes its connected monodromy $H$ normal in $G$; see
\cite[\S5, Theorem~1]{AndreMonodromy}.  In adjoint notation write
\[
                            G=H\times Q.
\]
The $Q$-period is constant on the lifted $Z$, by the fixed-part
theorem after killing finite monodromy.  The varying monodromy
domain on $Z$ is therefore the domain for $H$.  Write $J_H$ for
its line-orbit projection and $d_H=\dim J_H$.
In local flat coordinates write $j=(j_H,j_Q)$, with targets $J_H$
and $J_Q$, the products of line orbits for the two groups.  We will
show that the $H$-line map has full rank $d_H$ on $Z$, whereas the
$Q$-line map and the full $Q$-period map have equal generic rank
on the ambient base.

For the first assertion, put
$m=\dim Z$ and $d=\dim\check{\mathcal D}_H$, where
$\check{\mathcal D}_H$ is the compact dual of the connected-monodromy
period domain.  The projection
$\check{\mathcal D}_H\to J_H$ is a surjective morphism of
homogeneous projective flag varieties.  All its fibers have
dimension $d-d_H$.  Let
$V\subset Z\times\check{\mathcal D}_H$ be the algebraic locus
where the projected line equals its value on $T$.  It follows that
$\dim V=m+d-d_H$.  The period graph has dimension $m$ in an
ambient variety of dimension $m+d$, so its expected intersection
dimension with $V$ is $m-d_H$.

By \cite[Theorem~1.1]{BakkerTsimermanAxSchanuel}, an intersection
component of larger dimension projects into a proper algebraic weak
Mumford--Tate subvariety of $Z$.  This is impossible for the component
containing the lifted period graph over $T$, since $T$ is Zariski
dense in $Z$.  Conversely, a map
to $J_H$ has fibers of dimension at least $m-d_H$.  At the chosen
smooth point the restricted fiber is exactly $T$, because the
ambient local fiber is $T\subset Z$.  Equality follows, proving the
claim.  The cited theorem applies to the rationally polarized
integral adjoint variation on a smooth algebraic base and uses
precisely its connected-monodromy domain.

Let $q$ denote the full $Q$-period map and set
$s_Q=\rank(dq)$. Choose the smooth point $x\in T\cap Z_{\mathrm{reg}}$
inside the ambient constant-rank neighborhood fixed above. Since
$T$ is Zariski dense in $Z$, such a point exists. The constant-rank
theorem gives $\ker(dj_x)=T_xT\subset T_xZ$, and hence
\[
 \rank(dj_x|_{T_xZ})=\dim Z-\dim T=d_H.
\]
The $Q$-period is constant on $Z$, so $T_xZ\subseteq\ker(dq_x)$.
Thus, at this same point,
\[
 \ker(dj)\subseteq\ker(dq),\qquad
             dj_H(\ker(dq_x))=T_{j_H(x)}J_H.
\]
The first inclusion uses $T_xT\subset T_xZ$; the second is the
just-established rank calculation on $T_xZ$.
The differential $dj_Q$ factors through $dq$.  If $dj_Q(v)=0$,
choose $w\in\ker(dq)$ with $dj_H(w)=dj_H(v)$.  Then
$v-w\in\ker(dj)$, and therefore $dq(v)=0$.  We conclude
\begin{equation}
\label{period:rankdecomposition}
       \rank(dj_Q)=s_Q,\qquad r=d_H+s_Q.
\end{equation}
These are generic identities by our choice of the point.

The factor-line construction, after a common positive rescaling
and a further finite cover, now gives nef lines with
\begin{equation}
\label{period:linedecomposition}
       \pi^*P=P_H+P_Q,\qquad
       \nu(P_H)=d_H,
       \qquad P_Q=q^*B_Q.
\end{equation}
Here $q:\widetilde S\to S_Q$ is a resolved projective quotient
for the full $Q$-period, constructed by
\Cref{period:quotient}, and $\dim S_Q=s_Q$.  The highest-line
rank on $S_Q$ is $s_Q$ by
\eqref{period:rankdecomposition}.  Only the quotient construction,
not adjoint positivity, has been used to construct $S_Q$.

\smallskip\noindent
\emph{Reduction for a boundary component.}
For a fixed marked $D_i$, use a common finite Galois cover that
preserves the factors and makes boundary monodromy unipotent.
Resolve it and the map $q$ equivariantly, and let $E_i$ be the
strict transform of a component above $D_i$ that is generically
finite over $D_i$.  The inertia group at its generic point fixes
$E_i$ pointwise.  Further rescaling does not change any numerical
dimension.  Nefness gives
\[
                     \nu(P_H|_{E_i})\leq d_H.
\]
For example, every intersection involving $E_i$ is bounded by
the corresponding intersection with a sufficiently ample
effective divisor; the vanishing of higher powers of $P_H$
then gives this inequality.  If $E_i$ does not dominate $S_Q$,
\[
                     \nu(P_Q|_{E_i})\leq s_Q-1,
\]
and the binomial expansion for the two nef lines in
\eqref{period:linedecomposition} proves the strict drop.
We may therefore suppose $E_i$ dominates $S_Q$.  At its generic
point the $Q$-variation is pulled back from the interior of
$S_Q$, so its local unipotent logarithm vanishes.

\smallskip\noindent
\emph{Infinite local monodromy.}
If the original local monodromy is infinite, its logarithm at
unipotent level is a nonzero rational element
$N\in\mathfrak h$.  Every nonzero rational element of a rational
simple factor has nonzero projection to every conjugate complex
factor.  Each retained rational factor has a visible complex
factor.  Hence $N$ acts nontrivially on $J_H$.

Let $w$ be the pure weight of the representation supplying $P_H$,
and let $F^a=\C v$ be its limiting highest line at a general
point of $E_i$.  This line has a unique weight contribution,
say of weight $w+k$.  It is primitive and $k\geq0$: a
nonprimitive contribution, or a contribution below the middle
weight, would come by monodromy Lefschetz from a nonzero
filtration piece above $F^a$.
Moreover
\begin{equation}
\label{period:nilpotenttop}
                    N^{k+1}v=0,\qquad N^kv\ne0.
\end{equation}
To see that this statement holds for the actual vector rather
than only its weight-graded image, use the Deligne splitting
of the limiting mixed Hodge structure.  The one-dimensional
highest step occupies a single summand, and $N$ has degree
$(-1,-1)$ for this functorial splitting.  Primitive vanishing
on the graded piece therefore gives the first equality in
\eqref{period:nilpotenttop}; the Lefschetz isomorphism gives
the second.  For a complex summand the same argument is made
in its conjugate-paired real variation and restricted to the
summand.

In nilpotent-orbit coordinates $[v]\in J_H$.  Since this orbit
is closed in projective space and invariant under $\exp(tN)$,
\begin{equation}
\label{period:fixedlocus}
       [N^kv]=\lim_{t\to\infty}\exp(tN)[v]
                \in J_H\cap\mathbb P(\ker N).
\end{equation}
This fixed locus is a proper closed subset of the irreducible
variety $J_H$, because $N$ acts nontrivially on it.  Its dimension
is less than $d_H$.  Restrict to a simply connected open part of the boundary stratum
on which $k$ is constant.  In a flat trivialization, both $N$ and
$W(N)$ are constant.  Project $N^kv$ to
$\Gr^{W(N)}_{w-k}$ and apply the inverse of the fixed Lefschetz
isomorphism
\[
 N^k:\Gr^{W(N)}_{w+k}\longrightarrow\Gr^{W(N)}_{w-k}.
\]
This recovers the weight-graded limiting highest line.  Every actual
vector $N^kv$ lies in $W(N)_{w-k}$, and its image in that graded
piece is nonzero.  Thus the graded line map factors through the
rational map on
\[
 J_H\cap\mathbb P(\ker N)\cap\mathbb P(W(N)_{w-k})
\]
induced by this fixed linear quotient and inverse.  This locus is
contained in the proper fixed locus in \eqref{period:fixedlocus},
and the rational map is defined at every point under consideration.
Its rank is less than $d_H$.
By \Cref{period:boundary},
\[
                       \nu(P_H|_{E_i})<d_H.
\]
Together with $\nu(P_Q|_{E_i})\leq s_Q$ this proves the
required inequality in the infinite-monodromy case.

\smallskip\noindent
\emph{Finite nonidentity local monodromy.}
Let $\gamma$ be the original finite local monodromy.  The
disconnected-monodromy argument in \Cref{period:adjoint-reduction}
shows that $\gamma$ preserves the orbit span $W$.  It therefore acts
on the highest-line orbit $J$, permuting its nontrivial simple-factor
flag varieties. Equivariance preserves the kernel of
$G_{\C}\to\operatorname{Aut}(J)$, which is the product of the
invisible simple factors. Thus visible and invisible factors cannot
be exchanged. At unipotent level its logarithm is zero, so the variation extends
with pure limiting flags.  The inertia generator fixes the generic
boundary point; equivariance therefore makes these flags fixed by
$\gamma$.  Suppose that the
boundary line retained rank $r=d_H+s_Q$.  Since $E_i$ dominates $S_Q$,
its full $Q$-period and its $Q$-line map both have rank $s_Q$.
Along a generic fiber of this $Q$-period, the $H$-line map
must consequently have rank $d_H$.  Its coordinates fill an
open subset of $J_H$ while the full $Q$-flags remain fixed.

The fixed-point condition for $\gamma$ rules out permutations
involving any visible $H$-factor, including a permutation
with a factor whose coordinate is held fixed.  On each visible
$H$-factor of the highest-line orbit it forces the identity action.
The action of a simple adjoint group on a nontrivial flag variety is faithful;
equivariance then makes the induced automorphism of that
simple group the identity as well.  Rationality forces the
same conclusion on all conjugate factors.  Thus $\gamma$
acts identically on $H$ and preserves $Q$.

Its $Q$-action fixes every generic $Q$-period, by the domination
of $S_Q$ and analytic continuation from the resulting open set of
lifted periods. Apply \Cref{period:stabilizer} to the descended
$Q$-variation.  This action is the identity.  Hence $\gamma$
is the identity on the whole retained adjoint system,
contradicting the marking.

In both cases nef binomial expansion and
\eqref{period:rankdecomposition} give
$\nu(\pi^*P|_{E_i})<r$.  Numerical dimension of a nef line is
unchanged by a generically finite pullback, so this is
\eqref{period:strictdrop}.  If $r=0$, the local line fiber is the whole source germ, so its
Zariski closure is the whole source and $H=G$, $Q=1$.  The
Ax--Schanuel calculation gives $d_H=0$.  Every retained rational
factor contains a factor acting nontrivially on the line orbit;
hence there can be no retained factor.  Generic immersivity of the
full retained period map then forces $S$ to be a point, with no
marked component.
\end{proof}

\subsection{Removing a boundary by counting sections}

The period argument must pass from a big logarithmic canonical
divisor to a big adjoint divisor without its boundary. The following
lemma isolates the last step: sections lost on the boundary grow
more slowly in the nef direction than sections on the whole space.
For a nef rational divisor $J$ on a smooth projective $n$-fold, set
\[
 \nu(J)=\max\{q\in\{0,\ldots,n\}:J^q\cdot H^{n-q}>0\},
\]
where $H$ is ample. This is independent of $H$ and equals zero
on a point. The proof below fixes the nef coefficient before
taking the asymptotic section limit.

\begin{lemma}\label{period:subtract}
Let $Z$ be a smooth projective variety, $J$ a nef rational divisor,
$A$ a rational divisor, and
$D=\sum_{i=1}^s d_iE_i$ an effective rational divisor with smooth
prime components. Suppose $A+D$ is big and
\[
 \nu(J|_{E_i})<r:=\nu(J)
 \quad\text{for every component of }D.
\]
Then $A+tJ$ is big for every sufficiently large rational $t$.
\end{lemma}

\begin{proof}
If $r=0$, the condition forces $D=0$, so $A$ is already big;
adding a nef divisor preserves bigness. This also covers dimension
zero with the usual convention. Assume $r>0$ and put $n=\dim Z$.
Choose an ample rational divisor $H_1$ with
$A+D-H_1$ rationally effective. For fixed rational $t\geq0$
and sufficiently divisible $m$,
\[
 h^0(Z,m(A+D+tJ))\geq h^0(Z,m(H_1+tJ)).
\]
Since $H_1+tJ$ is ample, asymptotic Riemann--Roch makes the
leading coefficient on the right
$(H_1+tJ)^n/n!$. As a polynomial in $t$, this intersection has
degree $r$ and positive leading coefficient.

Remove $mD$ one component at a time by divisor exact sequences.
Every lost quotient is a line bundle on an $E_i$ of the form
\[
 \mathcal O_{E_i}\left(m(A+D+tJ)-\sum_h k_hE_h\right),
 \qquad 0\leq k_h\leq md_h.
\]
Choose an ample integral divisor $B_0$ so large that a fixed
denominator-clearing multiple of $B_0-A-D$ is globally generated.
For each $h$, choose an ample integral divisor $B_h$ such that
both $B_h$ and $B_h+E_h$ are globally generated. Put
$H_2=B_0+\sum_h d_hB_h$. The difference between
$m(H_2+tJ)$ and the displayed quotient divisor is
\[
 m(B_0-A-D)+
 \sum_h\bigl((md_h-k_h)B_h+k_h(B_h+E_h)\bigr).
\]
It is globally generated in these divisible degrees. A section
not vanishing identically on $E_i$ gives an injection into
$\mathcal O_{E_i}(m(H_2+tJ)|_{E_i})$. There are exactly
$md_i$ quotient terms on $E_i$. Hence
\[
 h^0(Z,m(A+tJ))
 \geq h^0(Z,m(H_1+tJ))
       -m\sum_i d_i h^0(E_i,m(H_2+tJ)|_{E_i}).
\]
For each fixed $t$, divide by $m^n/n!$ and let $m$ tend to
infinity through the specified multiples. All the comparison
divisors are ample, so the resulting lower bound is
\[
 (H_1+tJ)^n
     -n\sum_i d_i(H_2+tJ)^{n-1}\cdot E_i.
\]
Each restriction polynomial has degree at most
$\nu(J|_{E_i})\leq r-1$, because the restricted divisor is nef.
The positive degree-$r$ term of the first polynomial therefore
dominates the loss for every sufficiently large $t$. The section
growth of $A+tJ$ is then positive in dimension $n$, which is
exactly bigness.
\end{proof}

Taking $A=K_Z$ and $D$ reduced gives the required removal of the
period boundary once its restriction ranks have been established.
The weights also allow simultaneous removal of other effective
divisors satisfying the same rank inequality, for example a
divisorial ramification term on a generically finite cover.


\begin{proof}[Proof of \Cref{period:adjoint}]
Apply \Cref{period:quotient}.  If its target is a point there
is nothing further to prove.  Otherwise mark the divisor $D$
as above.  The adjoint variation extends across every unmarked
boundary component, including their intersections away from
$D$, because its local monodromies there are trivial.  Its
period differential is generically injective.  Consequently
\cite[Theorem~1.1]{BrunebarbeCadorel} gives
\[
                              K_S+D\text{ big}.
\]
If an initial choice omitted codimension-two loci, apply the
theorem on a log resolution and push the sections down; the
resulting logarithmic canonical divisor pushes forward to
$K_S+D$.  By \Cref{period:rankloss}, each marked component
has strictly smaller numerical rank.  Apply \Cref{period:subtract} with $A=K_S$ and the marked reduced
divisor $D$. It gives $K_S+aP$ big.  Together with the
rational line identity in \Cref{period:quotient}, this proves
\eqref{period:conclusion}.  Since $P$ is nef, the same bigness
holds for every larger rational coefficient $a$.
\end{proof}
